\section{Parsed validation-snapshot provenance and observable-support policy}

This gate identifies the locked object precisely.  The committed hash covers a
release-local parsed validation snapshot.  The archive PDF, XML, and CIF
validation payloads advance through their own byte-lock gates.

\subsection{Field-level provenance}

Every snapshot field used by the support audit carries the chain
\[
\text{archive payload}
\longrightarrow
\text{source-field locator}
\longrightarrow
\text{parsed snapshot field}
\longrightarrow
\text{snapshot SHA-256}.
\]
The current evidence ledger records PDF page, report section, table/row key, the
snapshot JSON path, extraction method, and source-payload lock state.  Native XML
XPath or validation-CIF row locators join after those archive payloads are
byte-locked.  A later equivalence audit regenerates the snapshot from the locked
machine-readable source and compares every parsed field.

\subsection{Observable-aware local support}

The active derived observable uses
\[
\text{atom selector}=C_\alpha,
\qquad
O_{ij}=\mathbf 1\!\left[d_{C_\alpha,C_\alpha}(i,j)\leq 8.0\,\text{\AA}\right]
\]
inside the frozen sequence-separation boundary.  Local evidence is therefore
classified by its relevance to the selected atom:
\[
\begin{array}{c|c}
\text{condition} & \text{support action}\\ \hline
\text{selected }C_\alpha\text{ missing or unusable} & \texttt{quality\_excluded}\\
\text{unresolved selected-atom alternate location} & \texttt{quality\_excluded}\\
\text{selected-atom/backbone geometry flag} & \texttt{quality\_ambiguous}\\
\text{side-chain-only geometry flag} & \texttt{quality\_ambiguous}\\
\text{missing local model-to-data support} & \texttt{quality\_ambiguous}
\end{array}
\]
Geometry flags and local model-to-data support remain separate evidence classes.
No RSRZ or RSCC threshold is chosen after inspecting the entry values; the rule
uses an explicitly named native validation classification when available and
maps missing support to ambiguity.

The rematerialized residue state is
\[
N_{\rm residue}^{\rm supported}=0,
\qquad
N_{\rm residue}^{\rm ambiguous}=46,
\qquad
N_{\rm residue}^{\rm excluded}=0.
\]
All 946 pair rows are therefore quality-ambiguous and retain
\texttt{abstain} as the effective target state.  The coordinate-model fixture
still preserves 114 derived contact bits and 832 derived noncontact bits.

\subsection{Legacy-entry and scored-accession policy}

The 1CRN lane is frozen as
\[
\texttt{measurement\_lineage\_coordinate\_model\_abstention\_fixture}.
\]
Its measurement lineage, coordinate model, and abstention mask remain useful,
but the first quality-supported derived-contact score requires a separately
selected accession.  Eligibility criteria are frozen before accession selection
and before any AOD agreement values are read.  The selection method is the
lexicographically first accession among all entries satisfying the declared
method, byte-lock, machine-readable validation, local-support, selected-atom
coverage, alternate-location, chain, size, resolution, payload-status, and
motif-policy criteria.

The target-mask gate remains
\[
N_{\rm quality\ supported\ pairs}>0,
\]
while the AOD comparison join additionally requires a declared
alignment/projection rule and nonzero in-scope coverage.  No target join,
residual, or score is opened here.

\subsection{Ledger pointers}

\begin{verbatim}
manual-2/data/protein/
  pdb_external_validation_snapshot_evidence_locators.csv
  pdb_external_validation_outlier_observable_policy.csv
  pdb_external_legacy_entry_policy.csv
  pdb_external_scored_accession_eligibility_rule.csv
  pdb_external_validation_snapshot_provenance_manifest.json
\end{verbatim}
